\subsection{Feedback type, stability, and the asymmetry of Type~I/II}
\label{sec:stability}

The opposite-side override experiments show why the two feedback types behave so
differently under distillation (Table~\ref{tab:stability}).

\paragraph{\texttt{force\_ii} is stable; \texttt{teacher\_sign} is high-ceiling but
unstable.}
Forcing the bounded Type~II on disagreement samples turns airquality, a weak teacher from
which the baseline blend barely helps, into a significant gain ($+0.175$, $t=2.33$) by
acting on samples the baseline ignores. Because Type~II can only erode clauses, the update
is self-limiting and never diverges. Routing the feedback type by the teacher's error sign
(\texttt{teacher\_sign}) instead gives the best single result on the good teacher ccpp
($+0.529$, $t=3.09$), but it diverges on airquality (mean $-17.4$; individual seeds at
$-5$, $-21$, $-80$; standard deviation $25.4$). The cause is the unbounded feedback type:
on a disagreement sample where the student over-predicts (error $>0$) but the teacher
under-predicts, \texttt{teacher\_sign} fires Type~I, which pushes an already-too-high
prediction higher, and because Type~I is unbounded this forms a positive-feedback loop with
no upper limit. \texttt{force\_ii} avoids the problem because it only ever fires the
lowering Type~II direction and never the raising Type~I direction.

\paragraph{Bounding either the type or the magnitude removes the divergence.}
Two fixes both eliminate the instability (airquality standard deviation drops from $25.4$ to
$\approx\!0.23$; Table~\ref{tab:stability}). Dropping the Type~I branch
(\texttt{teacher\_sign\_ii}) is safe but inert: it removes the divergence and most of the
benefit (airquality $\approx\!0$, energy still significantly harmful). Bounding the
magnitude to the teacher's frozen prediction error (\texttt{teacher\_sign\_tmag}) is better
on every axis: airquality $+0.107$ (positive), energy harm roughly halved ($-0.152$ versus
$-0.377$), and the lowest variance throughout. This is the first mechanism that makes a bad
teacher meaningfully less harmful rather than merely amplifying it.

\begin{table}[t]
\centering
\footnotesize
\caption{Stability fixes on airquality and energy, mean \texttt{impr} (std).}
\label{tab:stability}
\begin{tabular}{llrr}
\toprule
dataset & mode & mean \texttt{impr} & std \\
\midrule
\multirow{4}{*}{airquality}
  & uniform            & $+0.078$  & $0.19$ \\
  & teacher\_sign      & $-17.39$  & $25.41$ \\
  & teacher\_sign\_ii  & $-0.001$  & $0.23$ \\
  & teacher\_sign\_tmag& $\mathbf{+0.107}$ & $0.25$ \\
\midrule
\multirow{4}{*}{energy}
  & uniform            & $-0.062$  & $0.13$ \\
  & teacher\_sign      & $-0.377$  & $0.25$ \\
  & teacher\_sign\_ii  & $-0.366$  & $0.22$ \\
  & teacher\_sign\_tmag& $\mathbf{-0.152}$ & $0.11$ \\
\bottomrule
\end{tabular}
\end{table}

The broader lesson is that the TM's feedback asymmetry acts as a safety property: any
distillation mode confined to the bounded Type~II direction is inherently stable, whereas
modes that can fire the unbounded Type~I in the wrong direction require an explicit
magnitude bound.
